\documentstyle[%


righttag%


,11pt%


,amssymb%


,russian%               remove in Eng.


,izvru%                 Set izven% in Eng.


]{amsart}





%





\newcommand{\cp}{\operatorname{cap}}


\newcommand{\ACL}{\operatorname{ACL}}


\newcommand{\loc}{\operatorname{loc}}


\newcommand{\Lip}{\operatorname{Lip}}


\newcommand{\dist}{\operatorname{dist}}


\newcommand{\esssup}{\operatornamewithlimits{ess\,sup}}


\newcommand{\tts}[1]{}





\theoremstyle{definition}


\theorembodyfont{\rm}


\newtheorem{notenonum}{\hskip\parindent Замечание}


\renewcommand{\thenotenonum}{}





%\hoffset=-15mm%                  For Eng.


\setcounter{page}{11}


\god{2002}


\nomer{10~(485)} %                        Remove in Eng.


%\nomer{Vol.\,46, No.\,10}%               Set in Eng.


%\str{\pageref{begin}--\pageref{end}}%  Set in Eng.


\udk{517.548+517.988}





\author{\it C.K.\,Bодопьянов, А.Д.\,Ухлов}


% NB:   ^^ remove in Eng.


\title{Операторы суперпозиции в пространствах Соболева}


\thanks{Работа выполнена при частичной поддержке Российского фонда


фундаментальных исследований,


Совета государственной поддержки ведущих научных школ и INTAS-10170.}





\date{}


%\pagestyle{myheadings}%         Set in Eng.


\pagestyle{plain} %              Remove in Eng.


\begin{document}


%\thispagestyle{plain}%          Set in Eng.


\maketitle


\label{begin}











\section*{Введение}





Рассмотрение операторов суперпозиции в пространствах Соболева


связано с обобщением задачи Ю.Г.\,Решетняка об описании


всех изоморфизмов $\varphi^*$ однородных пространств Соболева


$L_n^1$, порожденных квазиконформными отображениями $\varphi$


евклидова пространства ${\Bbb R}^n$ по правилу


$\varphi^*(u)=u\circ\varphi$, поставленной в 1968~г.\ на первом


Донецком коллоквиуме по теории отображений. В [1] показано,


что таковыми являются структурные изоморфизмы пространств $L_n^1$


и только они. Предложенный в [1] подход к задаче Решетняка


естественно рассматривать в контексте предшествующих этому


результатов ([2], с.\,419--420). В теоремах Банаха,


Стоуна, Эйленберга, Аренса и Келли, Хьюита, Гельфанда и


Колмогорова получены условия на различные структуры пространства


непрерывных функций $C(S)$, изоморфизм которых определяет


топологическое пространство $S$ с точностью до гомеоморфизма.


Отметим здесь результат Стоуна, согласно которому $C(S)$ как


структурно упорядоченная группа определяет $S$. С другой стороны,


М.\,Накаи [3] и Л.\,Льюис [4] установили, что изоморфность алгебр


Ройдена равносильна квазиконформной эквивалентности


областей определения. Выделяя теперь в однородном пространстве


Соболева $L^1_n$ две структуры: векторной решетки и


полунормированного пространства --- получим ситуацию, в


алгебраическом смысле близкую к работе Стоуна, а в метрическом ---


к работе Накаи. Такой взгляд на задачу является наиболее


естественным, т.\,к. все еще дает возможность восстановить


отображение, несмотря на минимум ``материала'' для его нахождения,


доказать его непрерывность и установить его метрические свойства.





В рамках найденного в [1] подхода к проблеме Решетняка возникает


следующая задача: какие метрические и аналитические свойства имеет


измеримое отображение $\varphi$, индуцирующее изоморфизм


$\varphi^*$ по правилу $\varphi^*(f)=f\circ\varphi$, $f\in L_n^1$?


Варьируя функциональное пространство $L_n^1$, мы каждый раз


приходим к новой задаче: пространства Соболева $W^1_p$,


$p>n$, рассмотрены в [5], однородные пространства Бесова


$b^{l}_{p}({\Bbb R}^n)$, $n>1$, $lp=n$, --- в [6] при $p=n+1$ и в


[7] при $p>n+1$, пространства Соболева $W^1_p$, $n-1<p<n$, --- в


[8], пространства Соболева $W^1_p$, $1\leq p<n$ (и пространства


потенциалов) --- в [9], трехиндексные шкалы пространств


Никольского--Бесова и Лизоркина--Трибеля (и их анизотропные


аналоги) --- в [10]. В [11] к задаче замены переменной в


пространствах Соболева применена теория мультипликаторов. Вывод


из [5]--[10] состоит в том, что изоморфность оператора


$\varphi^*$ влечет (в зависимости от соотношения между


показателями гладкости, суммируемости и размерностью) свойствa


отображения быть квазиконформным или квазиизометрическим в метрике


области определения, адекватной геометрии


функционального пространства.





Аналитические и метрические свойства гомеоморфизмов, индуцирующих


{\it ограниченные операторы} пространств Соболева $L^1_p$,


изучались в [12]--[18]. В [16]--[18] получено аналитическое описание


таких гомеоморфизмов без каких-либо априорных


предположений. {\it Гомеоморфизм $\varphi:D \to


D^{\prime}$ между областями $D$ и $D^{\prime}$ в ${\Bbb R}^n$,


$n\geq2$, индуцирует ограниченный оператор


$\varphi^*:L^1_p(D^{\prime})\to L^1_p(D)$, $p\in[1,\infty)$, по


правилу $\varphi^*(f)=f\circ\varphi$ тогда и только тогда, когда


отображение $\varphi$ принадлежит $W^1_{1,\loc}(D)$ и


$|D\varphi(x)|^p\leq K_p|\det D\varphi(x)|$ почти всюду в $D$}.


Заметим, что при $p=n$ этот класс отображений совпадает с классом


квазиконформных отображений, а при $p=1$ такие отображения названы


В.Г.\,Мазьей в [12] субареальными. Известно также, что


квазиконформное отображение может быть определено через


метрические термины как гомеоморфизм, обладающий ограниченным


искажением (см. [19]--[21]). Аналог метрического определения при


$p\neq n$ для гомеоморфизмов, индуцирующих ограниченный


оператор $\varphi^*:L^1_p(D^{\prime})\to L^1_p(D)$ при


$n-1<p<\infty$ был получен в [22]. В [23] изучался класс


гомеоморфизмов, индуцирующих ограниченный оператор


$\varphi^*:L^1_p(D^{\prime})\to L^1_q(D)$ при $1\leq q<p<\infty$.


При $n-1<q<p=n$ этот класс гомеоморфизмов совпадает с классом


квазиконформных в среднем отображений, которые при некоторых


аналитических ограничениях изучались рядом авторов (напр.,


[24]). Некоторые применения квазиконформных в среднем отображений к


теории вложений функций классов Соболева имеются в [25].





В последнее время интенсивно развивается теория отображений с интегрируемым


искажением, обобщающая теорию отображений с ограниченным искажением (см.


[26], где приведена библиография по этому вопросу).





Основы теории $(p,q)$-квазиконформных гомеоморфизмов, являющихся


естественным обобщением квазиконформных отображений, заложены в [27].





Результаты данной работы можно рассматривать как


естественное развитие результатов работы [27]. Типичная задача


\S\,1 состоит в том, чтобы найти аналитические условия


на отображение $\varphi$, которое в отличие от [27] не


предполагается гомеоморфным, такие, что $\varphi$ индуцирует по


правилу замены переменной ограниченный оператор


$\varphi^*$ пространств Соболева. Результаты первого параграфа


развивают соответствующие результаты, полученные в [7], [18] для


пространств Соболева с одинаковой степенью суммируемости.


В \S\,2 установлены условия, при выполнении которых отображение,


индуцирующее ограниченный оператор пространств Соболева,


обладает ${\cal N}^{-1}$-свойством. В \S\,3 для гомеоморфных


отображений, индуцирующих ограниченный оператор


пространств Соболева, получены условия локализации нормы.








\section{Аналитические характеристики классов отображений}





Пусть $D$ и $D^{\prime}$ --- открытые множества в


евклидовом пространстве ${\Bbb R}^n$, $n\ge1$. Рассмотрим задачу о


суперпозиции отображения $\varphi:D\to D^{\prime}$ с


функциями из пространств $L^1_p(D^{\prime})$,


$1\leq p\leq\infty$.





Напомним, что пространство Соболева $W^1_q(D)$, $1\leq


q\leq\infty$, состоит из локально-суммируемых функций,


имеющих первые обобщенные производные, и таких, что конечна


величина


$$


\|f\mid W^1_q(D)\|= \|f\mid L_q(D)\|+\|\nabla f\mid


L_q(D)\|, \ \ \text{где}\ \ \nabla f=\bigg(\frac{\partial


f}{\partial x_1},\dotsc,\frac{\partial f}{\partial x_n}\bigg).


$$


Пространство Соболева $L^1_q(D)$ рассматривается с полунормой


$\|f\mid L^1_q(D)\|=\|\nabla f\mid L_q(D)\|$.





Функция $f:D\to {\Bbb R}$ {\it абсолютно непрерывна на прямой $l$},


имеющей непустое пересечение с $D$, если она абсолютно непрерывна


на любом отрезке этой прямой, заключенном в $D$. Функция $f:D\to


{\Bbb R}$ принадлежит классу $\ACL(D)$ ({\it абсолютно непрерывна


на почти всех прямых}), если она абсолютно непрерывна на почти


всех прямых, параллельных любой координатной оси. Отметим, что $f$


принадлежит пространству Соболева $L^1_{1}(D)$ тогда и только


тогда, когда $f$ локально-суммируема и ее можно изменить на


множестве нулевой меры так, что измененная функция принадлежит


$\ACL(D)$ и ее частные производные $\frac{\partial f}{\partial


x_i}(x)$, $i=1,\dotsc,n$, существующие почти всюду, суммируемы в


$D$. Заметим, что первые обобщенные производные функции $f$


совпадают почти всюду с обычными частными производными (напр., [12]).





Отображение $\varphi:D\to D^{\prime}$ принадлежит классу


$\ACL(D)$, если его координатные функции $\varphi_j$ можно


изменить на множестве нулевой меры так, чтобы они


принадлежали $\ACL(D)$, $j=1,\dotsc,n$. Отображение


$\varphi:D\to D^{\prime}$


принадлежит классу $W^1_{q,\loc}(D)$,


если оно принадлежит классу $\ACL(D)$ и его координатные функции


$\varphi_j$ принадлежат $W^1_{q}(V)$ для всякой компактной области


$V\subset D$, $j=1,\dotsc,n$.





Если отображение $\varphi:D\to D^{\prime}$ принадлежит классу


$\ACL(D)$, то для почти всех точек


открытого множества $D$ определены


формальная матрица Якоби $D\varphi(x)=\big(\frac{\partial


\varphi_i}{\partial x_j}(x)\big)$, $i,j=1,\dotsc,n$, и ее якобиан


$J(x,\varphi)=\det D\varphi(x)$. Норма $|D\varphi(x)|$ матрицы


есть норма линейного оператора, определяемого этой матрицей, в


евклидовом пространстве ${\Bbb R}^n$.





Будем говорить, что отображение


$\varphi:D\to D^{\prime}$


{\it порождает ограниченный оператор


вложения} $\varphi^{*} :L_{p}^{1}(D^{\prime})\to L_{q}^{1}(D)$,


$1\leq q\leq p \leq\infty$,


по правилу суперпозиции


$\varphi^{*}f=\ f\circ \varphi$, если существует постоянная


$K<\infty$ такая, что


$$


\|\varphi^{*}f\mid L_{q}^{1}(D)\|\leq


K\|f\mid L_{p}^{1}(D^{\prime})\|


\quad\text{для любой функции $f\in L_{p}^{1}(D^{\prime})$}.


$$


Аналогично определяется ограниченный оператор вложения в


пространствах Лебега и пространствах Соболева с основной нормой.





При описании отображений, порождающих по правилу суперпозиции


ограниченные операторы вложения пространств Соболева,


важную роль играют функции, определенные на системе


открытых подмножеств открытого множества $D$.





Напомним, что отображение $\Phi$,


определенное на открытых подмножествах из $D$


и принимающее неотрицательные значения, называется


{\it конечно $\vartheta$-квазиаддитивной} функцией множества [28], если


\begin{itemize}


\item[1)] для всякой точки $x\in D$ существует $\delta$,


$0<\delta<\dist (x,\partial D)$, такое, что


$0\leq\Phi(B(x,\delta))<\infty$ (здесь и далее


$B(x,\delta)=\{y\in{\Bbb R}^n:|y-x|<\delta\}$);


\item[2)] для всякого конечного набора $U_i\subset U\subset D$,


$i=1,2,\dotsc,k$, попарно непересекающихся открытых


множеств имеет место неравенство


$\sum\limits_{i=1}^{k}\Phi(U_i) \leq\vartheta \Phi(U)$.


\end{itemize}





Очевидно, неравенство во втором условии этого определения можно


распространить на счетную совокупность попарно


непересекающихся открытых множеств из~$D$, так что конечно


$\vartheta$-квазиаддитивная функция множества является также и


{\it счетно $\vartheta$-квазиаддитивной}.





Если вместо второго условия предположить, что


для всякого конечного набора $U_i\subset D$,


$i=1,2,\dotsc,k$, попарно непересекающихся открытых множеств


имеет место равенство


$$


\sum\limits_{i=1}^{k}\Phi(U_i)=


\Phi\Bigl(\bigcup\limits_{i=1}^{k}U_i\Bigr),


$$


то такая функция


множества называется {\it конечно-аддитивной}.


Если равенство в этом условии можно распространить на счетную 


совокупность попарно непересекающихся открытых множеств


из $D$, то такую функцию множества будем называть


{\it счетно-аддитивной}.





Отображение $\Phi$,


определенное на открытых подмножествах из $D$


и принимающее неотрицательные значения, называется {\it


монотонной} функцией множества [28], если


$\Phi(U_1)\leq\Phi(U_2)$ при условии, что


$U_1\subset U_2\subset D$ --- открытые множества.








{\it Bepхняя} и {\it нижняя} производные каждой из выше определенных


(квази)аддитивных функций множества, заданных на открытых


множествах из $D$, определяются следующим образом:


$$


\ov\Phi{}'(x)=\lim_{h\to0}\sup_{\delta<h}


\frac{\Phi(B_{\delta})}{|B_{\delta}|}\quad


\text{и}\quad


\underline\Phi'(x)=\lim_{h\to0}\inf_{\delta<h}


\frac{\Phi(B_{\delta})}{|B_{\delta}|},


$$


где точная верхняя и нижняя грани берутся по всем шарам


$B_{\delta}=B(y,\delta)\ni x$,


$B_{\delta}\subset D$,


радиус ${\delta}$ которых меньше $h$.


Заметим, что верхняя производная


$\ov\Phi{}'(x)$ и нижняя производная


$\underline\Phi'(x)$ являются


борелевскими функциями [28].





Сформулируем в нужной нам форме результат из [28].





\begin{prop}[{\series{m}\shape{n}\selectfont [28]}]


Пусть конечно $\vartheta$-квазиаддитивная функция


множества $\Phi$ определена на


открытых подмножествах открытого множества $D\subset{\Bbb R}^n$. Тогда


\begin{itemize}


\item[a)] для любого открытого множества $U\subset D$


справедливо неравенство


$$


\int_{U}\ov\Phi{}'(x)\,dx\leq \vartheta\Phi(U);


$$


\item[б)] для почти всех точек $x\in D$


верхняя производная конечна и


$$


\ov\Phi{}'(x)\leq \vartheta\underline\Phi'(x).


$$


\end{itemize}


Если $\vartheta=1$, то для почти всех точек $x\in D$


существует конечная производная


$$


\lim_{\delta\to 0,\,B_{\delta}\ni x}


\frac{\Phi(B_{\delta})}{|B_{\delta}|}=\Phi'(x).


$$


\end{prop}





Неотрицательная функция $\Phi$, определенная на некоторой


совокупности измеримых множеств из открытого множества $D$ и


принимающая конечные значения, называется {\it абсолютно


непрерывной}, если для любого числа $\varepsilon>0$ найдется число


$\delta>0$ такое, что $\Phi(A)<\varepsilon$ для всякого измеримого


множества $A\subset D$ из области определения, удовлетворяющего


условию $|A|<\delta$.





Напомним, что гомеоморфизм $\varphi$ открытых множеств $D$ и


$D^{\prime}$ называется $(p,q)$-ква\-зи\-кон\-форм\-ным [27],


если выполнено одно из следующих эквивалентных утверждений:





1)~отображение $\varphi$ индуцирует ограниченный оператор вложения


$$


\varphi^*:L^1_p(D^{\prime})\to L^1_q(D),\ \ 1\leq q\leq p<\infty,


$$


по правилу $\varphi^*(f)=f\circ\varphi$;





2)~отображение $\varphi$ принадлежит $W^1_{1,\loc}(D)$


и конечна величина


$$


K_{p,q}(D)=\|K_p(\cdot)\mid L_{\varkappa}(D)\|,


$$


где $K_p(x)=\inf\{k:


|D\varphi|(x)\leq k|J(x,\varphi)|^{1/p},\ x\in D \}$,


а число $\varkappa$ определяется из соотношения


$1/\varkappa = 1/q - 1/p$ при $q<p$


и $\varkappa=\infty$ при $q=p$,


причем норма оператора $\varphi^*:L^1_p(D^{\prime})\to L^1_q(D)$


эквивалентна $K_{p,q}(D)$.





При $p\in(1,\infty)$ утверждения 1, 2 эквивалентны утверждению





3)~для любого конденсатора $(F_0,F_1)\subset D'$


справедливы неравенства


\begin{alignat*}{2}


\cp_{p}(\varphi^{-1}(F_0),\varphi^{-1}(F_1);D)&\leq K^p


\cp_{p}(F_0,F_1;D^{\prime})&\ \ &\text{при }\ 1<q=p<\infty \\


%


\intertext{и}


\cp_{q}^{1/q}(\varphi^{-1}(F_0),\varphi^{-1}(F_1);D) &\leq


\Phi(D^{\prime}\setminus(F_0\cup F_1))^{\frac{1}{\varkappa}}


\cp_{p}^{1/p}(F_0,F_1;D^{\prime}) &\ \ &\text{при }\ 1< q<p<\infty,


\end{alignat*}


где $\Phi$ --- ограниченная монотонная


счетно-аддитивная функция, определенная на открытых подмножествах


открытого множества $D'$.





При $p\in(n-1,\infty)$ утверждения 1, 2, 3 эквивалентны утверждению





4)~существуют постоянная $K<\infty$ и при $q<p$ ограниченная


монотонная абсолютно непрерывная счетно-аддитивная функция


$\Psi$, определенная на открытых подмножествах открытого множества $D$,


такие, что


$$


\underset{r\to 0}{\varlimsup}{}


\frac{L_{\varphi}^{p}(x,r)r^{{n}-p}} {|\varphi(B(x,\lambda


r))|}\bigg/ \biggl(\frac{\Psi(B(x,\lambda


r))}{|B(x,r)|}\biggr)^{\frac{p-q}{q}}\leq K


$$


для всех точек


$x\in D$ (при $p=q$ знаменатель этого выражения, стоящий под косой


линией, равен~$1$). Здесь


$L_{\varphi}(x,r)=\max\limits_{|x-y|=r}|\varphi(x)-\varphi(y)|$


при условии, что $r<\dist(x,\partial D)$, $\lambda>1$ ---


фиксированное число.





Понятие $(p,q)$-квазиконформного отображения естественным образом


обобщает понятие квазиконформного отображения (при $p=q=n\geq2$


данное определение совпадает с определением квазиконформного


отображения), и $(p,q)$-квазиконформные отображения представляют


пример топологических отображений с интегрируемым


искажением [26].





Рассмотрим теперь суперпозицию функций, принадлежащих пространству


Соболева $L^1_p(D^{\prime})$, с отображением $\varphi : D\to


D^{\prime}$, $D$ и $D^{\prime}$ --- открытые множества в ${\Bbb R}^n$.





\begin{prop}


Пусть отображение $\varphi


: D\to D^{\prime}$ обладает свойством, что для всякой


функции $f\in L^1_p(D^{\prime})$, $1\leq p\leq\infty$,


суперпозиция $\varphi^{\ast}f = f\circ\varphi


\in\ACL(D)$. Тогда отображение $\varphi$ принадлежит


$\ACL(D)$.


\end{prop}





\begin{pf}


Рассмотрим финитные функции $\xi_N\in


C^{\infty}_0({\Bbb R}^n)$, $N\in {\Bbb N}$, равные единице на


компактном множестве $\ov{B(0,N)}$. Для любой координатной


функции $y_j : D^{\prime}\to {\Bbb R}$, $j=1,\dotsc,n$,


произведение $\xi_N(y)\cdot y_j \in L^1_p(D^{\prime})$. Тогда


функции $\varphi^{\ast}(\xi_N y_j)(x) =


\xi_N(\varphi(x))\varphi_j(x)$ принадлежат классу $L^1_q(D)\subset


\ACL(D)$.





Фиксируем произвольно координатную ось и рассмотрим семейство


$L_N$, состоящее из всех прямых $l$, параллельных выбранной


координатной оси и таких, что на каждой прямой $l\in L_N$ функция


$\varphi^{\ast}(\xi_N y_j)$, $N\in {\Bbb N}$, абсолютно непрерывна.


Так как множество $\{\varphi^{\ast}(\xi_N y_j)\}_{N\in {\Bbb N}}$


счетное, то проекция вдоль выбранной оси множества


$L=\bigcap\limits_{N=1}^\infty L_N$ на подпространство ${\Bbb


R}^{n-1}$ имеет полную $(n-1)$-мерную меру Лебега.





Теперь покажем, что отображение $\varphi$ абсолютно непрерывно на


любой прямой $l\in L$. Прежде всего заметим, что т.\,к.


$\xi_N(y)\cdot y_j = y_j$ для $y\in \ov{B(0,N)}$ и функции


$\varphi^{\ast}(\xi_N y_j)(x)$ непрерывны на прямой $l$, то


множество $\varphi^{-1}(B(0,N))\cap l$ открыто на прямой $l$. На


прямой $l$ выберем любые точки $x_1,x_2\subset D$ так, чтобы


отрезок $[x_1,x_2]\subset D$. Так как множество


$\varphi^{-1}(B(0,N))\cap l$ открыто на прямой $l$, и


$[x_1,x_2]\subset


\bigcup\limits_{N=1}^\infty\varphi^{-1}(B(0,N))$, то найдется


номер $N_1$ такой, что $[x_1,x_2]\subset\varphi^{-1}(B(0,N_1))$.


Следовательно, $\varphi$ абсолютно непрерывно на


$[x_1,x_2]$, т.\,к.


$$


\varphi|_{[x_1,x_2]}=(\varphi^{\ast}(\xi_{N_1}


y_1)|_{[x_1,x_2]}, \ldots, \varphi^{\ast}(\xi_{N_1}


y_n)|_{[x_1,x_2]}),


$$


а функции $\varphi^{\ast}(\xi_{N_1}


y_j)\in\ACL([x_1,x_2])$, $j=1,\dotsc,n$.


\end{pf}





В следующем утверждении описывается свойство аддитивности нормы оператора


суперпозиции, суженного на финитные функции.





\begin{lem}


Пусть отображение $\varphi :


D\to D^{\prime}$ порождает ограниченный оператор


вложения


$$


\varphi^{\ast} : L^1_p(D^{\prime})\to


L^1_q(D),\quad 1\leq q<p\leq\infty.


$$


Тогда


$$


\Phi(A^{\prime})=\sup_{f\in


\overset{\circ}{L}{}_{p}^{1}(A^{\prime})} \bigg(


\frac{\|\varphi^{\ast} f\mid {L}_{q}^{1}(D)\|}


{\|f\mid\overset{\circ}{L}{}_{p}^{1}(A^{\prime})\|}


\bigg)^{\sigma},\quad \text{где}\quad


\sigma=


\begin{cases}


\frac{pq}{p-q} & \text{при}\ p<\infty;\\


q\ & \text{при}\ p=\infty,


\end{cases}


$$


является ограниченной монотонной счетно-аддитивной функцией,


определенной на открытых ограниченных множествах


$A^{\prime}\subset D^{\prime}$.


\end{lem}





\begin{pf}


Очевидно, 


$\Phi(A_1^{\prime})\leq\Phi(A_2^{\prime})$,


если $A_1^{\prime}\subset A_2^{\prime}$.





Пусть


$A_{i}^{\prime}$, $i\in {\Bbb N}$, --- открытые попарно


непересекающиеся множества в $D^{\prime}$,


$A^{\prime}_0=\bigcup\limits_{i=1}^{\infty}A_i^{\prime}$.


Рассмотрим такую функцию $f_i\in\overset{\circ}{L}{}_{p}^{1}


(A_i^{\prime})$, чтобы одновременно выполнялись условия


$\|\varphi^{\ast} f_i\mid{L}_{q}^{1} (D)\|\geq


\big(\Phi(A_i^{\prime})(1-\frac{\varepsilon}{2^i})\big)^{\frac{1}{\sigma}}


\|f_i\mid\overset{\circ}{L}{}_{p}^{1}(A_i^{\prime})\|$


и


$\|f_i\mid\overset{\circ}{L}{}_{p}^{1}(A_i^{\prime})


\|^p=\Phi(A_i^{\prime})(1-\frac{\varepsilon}{2^i})$


при $p<\infty$


($\|f_i\mid\overset{\circ}{L}{}_{p}^{1}(A_i^{\prime})


\|=1$ при $p=\infty$), $\varepsilon\in(0,1)$. Полагая


$f_N=\sum\limits_{i=1}^{N}f_i$ и применяя неравенство Г\"ельдера


при $1\leq q<p\leq\infty$ (случай равенства), получим


\begin{multline*}


\|\varphi^{\ast} f_N\mid{L}_{q}^{1} (D)\|\geq


\bigg(\sum_{i=1}^{N}


\left(\Phi(A_i^{\prime})\left(1-\frac{\varepsilon}{2^i}


\right)\right)^{\frac{q}{\sigma}}


\|f_i\mid\overset{\circ}{L}{}_{p}^{1}(A_i^{\prime})


\|^q\bigg)^{1/q}= \\


=\bigg(\sum_{i=1}^{N}\Phi(A_i^{\prime})


\left(1-\frac{\varepsilon}{2^i}\right)\bigg)^{\frac{1}{\sigma}}


\bigg\|f_N\,\Big|\, \overset{\circ}{L}{}_{p}^{1}


\Big(\bigcup_{i=1}^{N}A_i^{\prime}\Big)\bigg\|


\geq \bigg(\sum_{i=1}^{N}\Phi(A_i^{\prime})


-\varepsilon\Phi(A_0^{\prime}) \bigg)^{\frac{1}{\sigma}}


\bigg\|f_N\,\Big|\, \overset{\circ}{L}{}_{p}^{1}


\Big(\bigcup_{i=1}^{N}A_i^{\prime}\Big)\bigg\|,


\end{multline*}


т.\,к. множества, на которых функции $\nabla \varphi^{\ast} f_i$


отличны от нуля, не пересекаются. Отсюда следует


$$


\Phi(A_0^{\prime})^{\frac{1}{\sigma}}\geq\sup\frac


{\|\varphi^{\ast} f_N\mid{L}_{q}^{1} (D)\|}


{\bigg\|f_N\,\Big|\, \overset{\circ}{L}{}_{p}^{1}


\Big(\bigcup\limits_{i=1}^{N}A_i^{\prime}\Big)\bigg\|}\geq


\bigg(\sum_{i=1}^{N}\Phi(A_i^{\prime})-\varepsilon\Phi(A_0^{\prime})


\bigg)^{\frac{1}{\sigma}},


$$


где точная верхняя грань берется по


всем функциям $f_N\in\overset{\circ}{L}{}_{p}^{1}


\Big(\bigcup\limits_{i=1}^{N}A_i^{\prime}\Big)$ указанного выше


вида. Так как $N$ и $\varepsilon$ произвольны, то


$$


\sum_{i=1}^{\infty}\Phi(A^{\prime}_i) \leq


\Phi\Big(\bigcup_{i=1}^{\infty}A^{\prime}_i\Big).


$$


Справедливость обратного неравенства проверяется непосредственно.


\end{pf}





Отображение $\varphi : D\to D^{\prime}$ {\it обладает $\cal N$-свойством}


Лузина, если образ всякого множества нулевой меры есть множество меры


нуль.





Для доказательства следующей теоремы нам понадобится





\begin{prop}[{\series{m}\shape{n}\selectfont [29], следствие 5.1}]


Предположим, что отображение $\varphi:A\to{\Bbb R}^n$,


$A\subset{\Bbb R}^n$ --- измеримое множество, имеет


аппроксимативные частные производные на множестве~$A$. Тогда


существует множество $\Sigma_{\varphi}\subset A$ нулевой меры


такое, что для любой неотрицательной измеримой функции $u:A\to{\Bbb


R}$ справедлива следующая формула замены переменной в интеграле


Лебега\/${:}$


\begin{equation}


\int_Au(x)|J(x,\varphi)|\,dx= \int_{{\Bbb


R^n}}\bigg(\, \sum\limits_{x\in


\varphi^{-1}(y)\cap(A\setminus\Sigma_{\varphi})}u(x)\bigg)\,dy.


\end{equation}


Если $\varphi$ обладает $\cal N$-свойством Лузина, то


$\Sigma_{\varphi}=\emptyset$.


\end{prop}





Заметим, что всякое отображение $\varphi\in W^1_{q,\loc}(D)$


обладает $\cal N$-свойством Лузина при $q>n$ (напр., [29]).





Отображение $\varphi : D\to D^{\prime}$ класса $\ACL(D)$


имеет {\it конечное искажение}, если


$D\varphi(x)=0$ почти всюду на множестве


$Z=\{x\in D:J(x,\varphi)=0\}$.





Для отображения $\varphi : D\to D^{\prime}$ класса $\ACL(D)$


определим функцию


\begin{equation}


D'\ni y\mapsto H_q(y)=


\begin{cases}


\bigg(\,\sum\limits\begin{Sb} x\in\varphi^{-1}(y)\setminus


\Sigma_{\varphi},\\ J(x,\varphi)\ne 0 \end{Sb}


\frac{|D\varphi|^q(x)}{|J(x,\varphi)|}\bigg)^{\frac{1}{q}}, &\ \\


0, & \text{если}\ \{


x\in\varphi^{-1}(y)\setminus \Sigma_{\varphi}:J(x,\varphi)\ne


0\}=\emptyset.


\end{cases}


\end{equation}


(Здесь и далее $\Sigma_{\varphi}\subset D$~--- множество из предложения~3.)





Две функции $f:D'\to{\Bbb R}$ и $g:D'\to{\Bbb R}$ называются


эквивалентными ($f\sim g$), если $\alpha f(x)\leq g(x)\leq\beta f(x)$


для всех $x\in D'$, $0<\alpha\leq\beta<\infty$~--- постоянные числа.





Следующее утверждение сформулировано в [30] для ограниченной


области $D^{\prime}$ (при $p=q$ в эквивалентной форме оно установлено


в [7], [18]). Для гомеоморфного отображения $\varphi$ оно


доказано в [16]--[18] при $p=q$ и в [23] при $q<p$.





\begin{thm}


Отображение $\varphi : D\to


D^{\prime}$ порождает ограниченный оператор вложения


$$


\varphi^{\ast} :


L^1_p(D^{\prime})\to L^1_q(D),\quad 1\leq q\leq p<\infty,


$$


тогда и только тогда, когда $\varphi$ принадлежит $\ACL(D)$, имеет


конечное искажение, и


$$


H_q(\cdot)\in


L_{\varkappa}(D'),


\quad\text{где}\quad \varkappa=


\begin{cases} \frac{pq}{p-q} & \text{при}\ q<p<\infty;\\


\infty & \text{при}\ q=p<\infty.


\end{cases}


$$


Норма оператора $\varphi^{*}


:L_{p}^{1}(D^{\prime})\to L_{q}^{1}(D)$ эквивалентна


$H_{p,q}(D^{\prime})=\|H_q(\cdot)\mid L_{\varkappa}(D')\|{:}$


$$


\alpha_{p,q}H_{p,q}(D^{\prime})\leq\|\varphi^{*}\|\leq


H_{p,q}(D^{\prime}),


$$


где $\alpha_{p,q}$ --- положительная постоянная.





Кроме того, если $V\subset D$ --- открытое множество и


$\varphi(V)$ содержится в открытом множестве $W$, то оператор


$$


\varphi_V^{\ast} : L^1_p(W)\to L^1_q(V),\quad


\varphi_V^{\ast}(f)=f\circ\varphi|_V, \quad 1\leq q\leq


p<\infty,


$$


также ограничен, и $\|\varphi_V^{\ast}\|\leq


\|H_q(\cdot)\mid L_{\varkappa}(W)\|$.


\end{thm}





Для доказательства теоремы при $q<p$ понадобится


представляющая независимый интерес





\begin{thm}


Отображение $\varphi : D\to


D^{\prime}$ порождает ограниченный оператор вложения $\varphi^{\ast} :


L^1_p(D^{\prime})\to L^1_q(D)$, $1\leq q< p<\infty$, тогда и


только тогда, когда $\varphi$ принадлежит $\ACL(D)$, имеет


конечное искажение и существует ограниченная монотонная


счетно-аддитивная функция $\Phi$, определенная на открытых


ограниченных подмножествах из $D^{\prime}$, такая, что


выполняются соотношения


$$


\alpha^{\varkappa}_{p,q}H^{\varkappa}_q(y)\leq\Phi^{\prime}(y)\leq


H^{\varkappa}_q(y)


$$


для почти всех $y\in D'$, $1/\varkappa = 1/q -


1/p$, гдe $\alpha_{p,q}$ --- постоянная из теоремы $1$.


\end{thm}





Для доказательства этих теорем потребуется оценка значений


монотонной счетно-аддитивной функции через кратность покрытия.





С помощью теоремы Безиковича [31] доказывается





\begin{lem}


Для всякого открытого множества


$U\subset {\Bbb R}^n$, $U\ne{\Bbb R}^n$, существует счетное


семейство шаров ${\cal B}=\{B_j\}$ такое, что


\begin{itemize}


\item[1)] $\bigcup\limits_{j}B_j=U;$


\item[2)] если $B_j=B_j(x_j,r_j)\in\cal B$,


то $\dist(x_j,\partial U)=12r_j;$


\item[3)] семейства ${\cal B}=\{B_j\}$ и $2{\cal B}=\{2B_j\}$, где


символ $2B$ обозначает шар удвоенного радиуса с центром в той же


точке, образуют конечнократное покрытие множества $U;$


\item[4)] если $2B_j=B_j(x_j,2r_j)$, $j=1,2$, --- пересекающиеся шары, то


$\frac{5}{7}r_1\leq r_2\leq\frac{7}{5}r_1;$


\item[5)] семейство $\{2B_j\}$ можно разбить на конечное число, 


зависящее только от размерности $n$,


наборов так, что внутри каждого набора шары не пересекаются.


\end{itemize}


\end{lem}





\begin{lem}


Пусть монотонная


счетно-аддитивная функция $\Phi$ определена на открытых


подмножествах открытого множества $D\subset {\Bbb R}^n$. Тогда для всякого


открытого множества $U\subset D$, $U\ne{\Bbb R}^n$,


существует последовательность шаров $\{B_j\}$ такая, что


\begin{itemize}


\item[1)] семейства $\{B_j\}$ и $\{2B_j\}$ образуют


конечнократное покрытие множества $U;$


\item[2)] $\sum\limits_{j=1}^{\infty}\Phi(2B_j)\leq \zeta_n\Phi(U)$,


где постоянная $\zeta_n$ зависит только от размерности $n$.


\end{itemize}


\end{lem}





{\bf Доказательство.}


В соответствии с леммой 2 построим последовательности


шаров $\{B_j\}$, $\{2B_j\}$ и


разобьем последовательность $\{2B_j\}$ на $\zeta_n$ подсемейств


$\{2B_{1j}\}_{j=1}^{\infty}, \dotsc, \{2B_{\zeta_nj}\}_{j=1}^{\infty}$


таким образом, что внутри каждого набора шары не пересекаются,


$2B_{ki}\cap 2B_{kj}=\emptyset$, если $i\ne j$, $k=1,\dotsc,\zeta_n$.


Следовательно,


$$


\sum_{j=1}^{\infty}\Phi(2B_j)=


\sum_{k=1}^{\zeta_n}\sum\limits_{j=1}^{\infty}\Phi(2B_{kj})\leq


\sum_{k=1}^{\zeta_n}\Phi(U)={\zeta_n}\Phi(U).\quad\square


$$





{\bf Доказательство теорем 1 и 2. Необходимость.} 1.


Принадлежность отображения $\varphi$ классу $\ACL(D)$ следует


из предложения 2. Докажем, что отображение $\varphi$ имеет


конечное искажение. По лемме~1 для любой функции $f\in


\overset{\circ}{L}{}_{p}^{1}(A')$ при $q\leq p$ выполняется


неравенство $\|\varphi^{\ast} f\mid{L}_{q}^{1}(D)\|


\leq\Phi(A')^{1/\sigma}


\|f\mid\overset{\circ}{L}{}_{p}^{1}(A')\|$, где


$A'\subset D^{\prime}$~--- открытое ограниченное подмножество


(при $q=p$ полагаем $\Phi(A)^{\frac{1}{\sigma}}=\|\varphi^*\|$).


Фиксируем срезку $\eta\in C_{0}^{\infty}({\Bbb R}^n)$, равную


единице на $B(0,1)$ и нулю вне шара $B(0,2)$. Подставляя в это


неравенство функции


$h_{j}(z)=(z-y)_{j}\eta(\frac{z-y}{r})$, $j=1,\dotsc,n$,


где $(z-y)_{j}$ обозначает $j$-ю координату вектора


$z-y$, $B(y,2r)\subset D^{\prime}$,


приходим к неравенству


\begin{equation}


\bigg(\,\int\limits_{\varphi^{-1}(B(y,r))}|D\varphi|^q(x)\,dx


\bigg)^{1/q}\leq C\Phi(B(y,2r))^{1/\sigma}|B(y,2r)|^{1/p},


\end{equation}


где $C$ --- некоторая постоянная, зависящая только от размерности $n$


и показателя $p$.





Пусть $Z=\{x\in D\mid J(x,\varphi)=0\}$. Покажем, что


\begin{equation}


\int_{Z}|D\varphi|^q(x)\,dx=0.


\end{equation}


По формуле (1) $|\varphi(Z\setminus \Sigma_{\varphi})|=0$.


Фиксируем число $\varepsilon>0$ и открытое множество $U$ такие,


что $U\supset \varphi(Z\setminus \Sigma_{\varphi})$ и


$|U|<\varepsilon$. Выберем согласно лемме 2 конечнократное


покрытие $\{B(y_i,r_i)\}$ открытого множества $U$ шарами такое,


что шары $B(y_i,2r_i)\subset U$, $i\in{\Bbb N}$, также образуют


конечнократное покрытие множества $U$, кратность $M$ покрытия


$\{B(y_i,2r_i)\}$ не зависит от множества $U$ и $\sum\limits_i


|B(y_i,2r_i)|<M\varepsilon$. Тогда, применяя неравенство (3) и


лемму 3, выводим


\begin{multline*}


\int\limits_{Z}|D\varphi|^q(x)\,dx= \int\limits_{Z\setminus


\Sigma_{\varphi}}|D\varphi|^q(x)\,dx\leq


\sum_{i=1}^{\infty}\;\int\limits_{\varphi^{-1}(B(y_i,r_i))}


|D\varphi|^q(x)\,dx \le \\


\leq


\begin{cases}


C^p\|\varphi^*\|^p\sum\limits_{i=1}^{\infty}|B(y_i,2r_i)|\leq


C_1^p\|\varphi^*\|^p|U| &\text{при}\ q=p;\\


C^q\sum\limits_{i=1}^{\infty}\Phi(B(y_i,2r_i))^{\frac{q}{\sigma}}


\bigl(|B(y_i,2r_i)|\bigr)^{\frac{q}{p}} \leq


C_2^q\Phi(U)^{\frac{p-q}{p}} |U|^{\frac{q}{p}}


& \text{при}\ q<p.


\end{cases}


\end{multline*}





Так как $\Phi(U)\leq \|\varphi^{\ast}\|^{\sigma}<\infty$ и


$\varepsilon>0$ --- произвольное число, то (4) доказано, и,


cледовательно, $|D\varphi|=0$ почти всюду на множестве $Z$.





2. Перейдем к оценке нормы. Подставляя в неравенство


$\|\varphi^{\ast} f\mid{L}_{q}^{1}(D)\|


\leq\Phi(A')^{1/\sigma}


\|f\mid\overset{\circ}{L}{}_{p}^{1}(A')\|$ функции


$h_{j}(z)=(z-y)_{j}\eta(\frac{z-y}{r})$, $j=1,\dotsc,n$,


$B(y,2r)\subset D^{\prime}$, приходим к неравенству


$$


\bigg(\,\int\limits_{\varphi^{-1}(B(y,r))} |D\varphi|^q\,dx


\bigg)^{1/q} \leq


\begin{cases}


C\Phi(B(y,2r))^{1/\sigma}|B(y,r)|^{1/p} &\text{при}\ q<p;\\


C\|\varphi^*\|\cdot|B(y,r)|^{1/p} &\text{при}\ q=p,


\end{cases}


$$


где $C$ --- некоторая постоянная.





При $q=p$ применим к левой части полученной оценки формулу (1).


Тогда из теоремы Лебега о дифференцировании интеграла вытекает,


что


$$


\bigg(\sum\begin{Sb} x\in\varphi^{-1}(y)\setminus


\Sigma_{\varphi},\\ J(x,\varphi)\ne0 \end{Sb}


\frac{|D\varphi|^p(x)}{|J(x,\varphi)|}\bigg)^{\frac{1}{p}} \leq


C\|\varphi^*\|


$$


для почти всех $y\in D^{\prime}$ таких, что


$\varphi^{-1}(y)\setminus (\Sigma_{\varphi}\cup Z)\ne\emptyset$,


и $J(x,\varphi)\ne0$,


если $x\in \varphi^{-1}(y)\setminus (\Sigma_{\varphi}\cup Z)$.


Отсюда вытекает оценка


$$


H_{p,p}(D^{\prime})=\|H_p(\cdot)\mid L_{\infty}(D')\|\leq C\|\varphi^*\|.


$$





В случае $q<p$ приведем полученное неравенство к виду


$$


\int\limits_{\varphi^{-1}(B(y,r))}|D\varphi|^q(x)\,dx \leq


\wt{C}^q\bigg(\frac{\Phi(B(y,2r))}{|B(y,2r)|}


\bigg)^{\frac{q}{\sigma}}|B(y,r)|


$$


и применим к левой части этого соотношения формулу (1) замены


переменной


\begin{multline*}


\int\limits_{\varphi^{-1}(B(y,r))} |D\varphi|^q(x)\,dx=


\int\limits_{\varphi^{-1}(B(y,r))\setminus Z}


|D\varphi|^q(x)\,dx =\\


=\int\limits_{B(y,r)}


\sum_{x\in\varphi^{-1}(y)\setminus(\Sigma_{\varphi}\cup Z)}


\frac{|D\varphi|^q(x)}{|J(x,\varphi)|} \,dy \leq


\wt{C}^q\bigg(\frac{\Phi(B(y,2r))}{|B(y,2r)|}


\bigg)^{\frac{q}{\sigma}}|B(y,r)|.


\end{multline*}


Из теоремы Лебега о дифференцировании интеграла и свойств


производной аддитивной функции множества (предложение 1)


вытекает


$$


H^{\varkappa}_q(y)=\bigg(\,\sum_{x\in\varphi^{-1}(y)\setminus


(\Sigma_{\varphi}\cup Z)}


\frac{|D\varphi|^q(x)}{|J(x,\varphi)|}\bigg)^{\frac{p}{p-q}} \le


\wt{C}^{\varkappa}\Phi'(y)


$$


для почти всех $y\in D^{\prime}$


таких, что $\varphi^{-1}(y)\setminus (\Sigma_{\varphi}\cup


Z)\ne\emptyset$, где $H_q(\cdot)$ определена формулой (2).





Интегрируя последнее неравенство по ограниченному открытому


множеству $V\subset D^{\prime}$, получим


\begin{multline*}


\|H_q(\cdot)\mid L_{\varkappa}(V)\|^{\varkappa}=\int_{V}


\bigg(\,\sum_{x\in\varphi^{-1}(y)\setminus


(\Sigma_{\varphi}\cup Z)}


\frac{|D\varphi|^q(x)}{|J(x,\varphi)|}\bigg)^{\frac{p}{p-q}} \,dy\le\\


\le \wt{C}^{\varkappa}\int_{V}\Phi'(y)\,dy \le


\wt{C}^{\varkappa}\Phi(V)\le


\wt{C}^{\varkappa}\|\varphi^{\ast}\|^\varkappa.


\end{multline*}


Так как выбор $V\subset D^{\prime}$


произволен, то $\|H_q(\cdot)\mid L_{\varkappa}(D^\prime)\|\leq


\wt{C}\|\varphi^{\ast}\|$.


\medskip





{\bf Достаточность}. Покажем, что для любой функции $f\in


L_p^1(D^{\prime})\cap C^1(D^{\prime})$ выполняется


неравенство $\|\varphi^*f\mid L_q^1(D)\|\le


H_{p,q}(D^{\prime})\|f\mid L_p^1(D^{\prime})\|$, $q\leq


p$. Так как $f\circ\varphi$ принадлежит классу $\ACL(D)$, то


\begin{multline*}


\|\varphi^* f\mid L_q^1(D)\|\le


\bigg(\int_{D\setminus Z}(|\nabla


f|(\varphi(x))|D\varphi|(x))^q\,dx\bigg)^{1/q}\le \\ \leq


\bigg(\int_{D'}|\nabla f|^q(y)


\bigg(\,\sum_{x\in\varphi^{-1}(y),\;x\notin


\Sigma_{\varphi}\cup Z}


\frac{|D\varphi|^q(x)}{|J(x,\varphi)|}\bigg) \,dy\bigg)^{\frac{1}{q}}.


\end{multline*}


Используя неравенство Г\"ельдера при $q<p$, выводим


оценку


\begin{equation}


\|\varphi^*f\mid L_q^1(D)\|\le


\bigg(\int_{D'}H^{\varkappa}_q(y)\,dy\bigg)^{1/\varkappa}


\bigg(\int_{D'}|\nabla f|^p(y) \,dy\bigg)^{1/p}


\end{equation}


(при $q=p$ левый сомножитель равен $H_{p,p}(D^{\prime})$). Отсюда вытекает


оценка для нормы


$$


\|\varphi^{\ast}\|\leq H_{p,q}(D^{\prime}).


$$





Чтобы получить поточечную оценку теоремы 2, заметим, что


неравенство (5) справедливо для любого открытого


ограниченного множества $A'\subset D'$ вместо $D'$ и


произвольной функции $f\in C_0^\infty(A^{\prime})$. Поэтому


$$


\Phi(A')\leq\int_{A'}H^{\varkappa}_q(y)\,dy


$$


для любого


открытого множества $A'\subset D'$. Дифференцируя левую и правую


части этого неравенства, выводим $\Phi'(y)\leq H^{\varkappa}_q(y)$


для почти всех $y\in D'$.





Чтобы распространить (5) на все функции $f\in


L_p^1(D^{\prime})$, $1<q\leq p<\infty$, аппроксимируем


$f$ последовательностью гладких функций $f_n\in L_p^1(D^{\prime})$


так, чтобы одновременно $\|f-f_n\mid


L_p^1(D^{\prime})\|\to0$ и $(f-f_n)\to0$ квазивсюду в


$D^{\prime}$ при $n\to\infty$. Поскольку прообраз


$\varphi^{-1}(S)$ множества $S\subset D^{\prime}$ нулевой емкости


имеет нулевую емкость, то


$\varphi^{\ast}(f_n)\to\varphi^{\ast}(f)$ квазивсюду в $D$


(определение и свойства емкости см., напр., в [12]). Из


этого наблюдения вытекает следующий вывод: распространение


по непрерывности оператора $\varphi^{\ast}$ с подпространства


$f\in L_p^1(D^{\prime})\cap C^\infty(D^{\prime})$ на $f\in


L_p^1(D^{\prime})$ совпадает с оператором подстановки


$\varphi^{\ast}$, $\varphi^{\ast}(f)=f\circ\varphi$.





Если $1=q<p$, то прием, описанный в предыдущем случае, позволяет


распространить этот оператор на все функции рассматриваемого


класса Соболева, единственное отличие в рассуждениях состоит в


том, что из последовательности $\varphi^{\ast}(f_n)$ при


условии, что $f_n$ сходится квазивсюду в $L_p^1(D^{\prime})$,


можно извлечь подпоследовательность, сходящуюся почти всюду в $D$.





Если же $q=p=1$, то вместо емкостной следует использовать более


грубую характеристику сходимости; если последовательность $f_n\in


L_1^1(D^{\prime})$ сходится к $f\in L_1^1(D^{\prime})$ в


$L_1^1(D^{\prime})$, то некоторая ее подпоследовательность


сходится почти всюду. Чтобы закончить доказательство, достаточно


воспользоваться следующим свойством: прообраз множества нулевой


меры есть множество нулевой меры при отображении $\varphi:D\to


D^{\prime}$, индуцирующем ограниченный оператор


$\varphi^{\ast}:L_1^1(D^{\prime})\to L_1^1(D)$ (см. ниже теорему 3).





Доказательство последнего утверждения вытекает из


$\varphi\in\ACL(D)$ и оценки $\wt H_q(y)\leq H_q(y)$ для


$y\in W$, где


$$


W\ni y\mapsto \wt H_q(y)=


\begin{cases}


\bigg(\,\sum\limits\begin{Sb} x\in(\varphi^{-1}(y)\cap V)\setminus


\Sigma_{\varphi},\\ J(x,\varphi)\ne 0\end{Sb}


\frac{|D\varphi|^q(x)}{|J(x,\varphi)|}


\bigg)^{\frac{1}{q}},\\


0, & \text{если}\ \{x\in\varphi^{-1}(y)\setminus


\Sigma_{\varphi}:J(x,\varphi)\ne0\}=\emptyset.


\end{cases}


$$





Монотонную функцию множества $\Phi$, определенную первоначально


лишь на открытых подмножествах открытого множества $D$,


распространим на измеримые множества $A\subset D$ по формуле


$$


\wt{\Phi}(A)=\inf_{U\supset A}\Phi(U).


$$


Очевидно, что условие монотонности обеспечивает


равенство $\wt{\Phi}(U)=\Phi(U)$ для любого открытого


множества $U\subset D$.$\quad\square$





\begin{cor}


Аддитивная функция множества


из леммы $1$ $U\mapsto\Phi(U)$, $U\subset D^{\prime}$ --- открытое


множество, абсолютно непрерывна. Кроме того, ee


распространение $\wt{\Phi}$ на измеримые множества


также абсолютно непрерывно.


\end{cor}





\begin{pf}


Очевидно, что оценка (5) справедлива не


только для открытого множества $D^{\prime}$, но и для


любого открытого ограниченного множества $U\subset


D^{\prime}$ и функции $f\in {\stackrel{\circ}{L}}{}_p^1(U)$.


Следовательно,


$$


\|\varphi^*f\mid L_q^1(D)\|\le


\bigg(\int_{U} H^{\varkappa}_q(y)\,dx\bigg)^{1/\varkappa}


\|f\mid\overset{\circ}{L}{}_p^1(U)\|.


$$


Отсюда в силу неравенства $\Phi(U)\leq \|H_q(\cdot)\mid


L_{\varkappa}(U)\|^{\varkappa}$ получаем абсолютную непрерывность


функции $\Phi$ и, следовательно, абсолютную непрерывность ее


продолжения~$\wt{\Phi}$.


\end{pf}





Пусть отображение $\varphi :D\to D^{\prime}$


порождает ограниченный оператор вложения


$\varphi^{*} :L_{p}^{1}(D^{\prime})\to L_{q}^{1}(D)$,


$1\leq q<p<\infty$. Для любого открытого


множества $W\subset D^{\prime}$ и функции $f\in L^1_p(W)\cap


C^1(W)$ суперпозиция $\varphi_W^*f=f\circ\varphi$ определена на


множестве $\varphi^{-1}(W)$. Множество $\varphi^{-1}(W)$ может


быть и не открытым, но оно обладает следующим свойством: для почти всех


прямых $l$ параллельных какой-нибудь координатной оси, например


$x_i$, $i=1,\dotsc,n$, пересечение $\varphi^{-1}(W)\cap l$ ---


открытое множество на прямой $l$, и суперпозиция $\varphi^*f$


абсолютно непрерывна на этом пересечении.


Следовательно, определена производная


$\frac{\partial \varphi_W^*f}{\partial x_i}(x)$


для почти всех $x\in \varphi^{-1}(W)\cap l$.


Таким образом, для почти всех $x\in\varphi^{-1}(W)$ определен градиент


$\nabla (\varphi_W^*f)(x)$, и полунорма


$$


\|\varphi_W^*f\mid L^1_q(\varphi^{-1}(W))\|\overset{\mathrm{def}}{=}{}


\bigg(\,\int\limits_{\varphi^{-1}(W)}|\nabla


\varphi_W^*f|^q(x)\,dx\bigg)^{\frac{1}{q}}.


$$


Аналогично соотношению (5) получим оценку


$$


\|\varphi_W^*f\mid L_q^1(\varphi^{-1}(W))\|\le


\bigg(\int_{W}H^{\varkappa}_q(y)\,dy\bigg)^{1/\varkappa}


\bigg(\int_{W}|\nabla f|^p(y) \,dy\bigg)^{1/p},


$$


из которой вытекает ограниченность оператора


$\varphi_W^*:L^1_p(W)\to L^1_q(\varphi^{-1}(W))$.





Определим функцию множества $W\mapsto[0,\infty)$ на открытых


множествах $W\subset D^{\prime}$ по следующей формуле:


$$


\Xi(W)=\sup_{f\in{L}_{p}^{1}(W)\cap C^1(W)}


\bigg(\frac{\|\varphi_W^{\ast} f\mid


{L}_{q}^{1}(\varphi^{-1}(W))\|} {\|f\mid


{L}_{p}^{1}(W)\|} \bigg)^{\varkappa}.


$$





Из теорем 1 и 2, следствия 1 и предложения 1 получим





\begin{cor}


Пусть $D$ и $D^{\prime}$ ---


открытые множества в ${\Bbb R}^n$. Отображение $\varphi :D\to D^{\prime}$


порождает ограниченный оператор вложения $\varphi^{*}


:L_{p}^{1}(D^{\prime})\to L_{q}^{1}(D)$, $1\leq q<p<\infty$, тогда


и только тогда, когда $\varphi$ принадлежит $\ACL(D)$, имеет


конечное искажение и существует ограниченная абсолютно непрерывная


счетно $\vartheta$-квазиаддитивная функция $\Xi$, определенная


на открытых множествах $W\subset D^{\prime}$, такая, что ее


верхняя производная $\ov\Xi{}^{\prime}(x)$ эквивалентна


$H^\varkappa_q(y){:}$


$$


\lambda_1 H^\varkappa_q(y)\leq


\ov\Xi{}^{\prime}(y)\leq H^\varkappa_q(y)


$$


для почти всех $y\in


D^{\prime}$, где $\lambda_1$ --- некоторая положительная


постоянная.





Норма оператора $\varphi^{*}_W$ эквивалентна


$\Big(\int\limits_W\ov\Xi{}^{\prime}(y)\,dy\Big)^{\frac{1}{\varkappa}}$


для любого открытого множества $W\subset D^{\prime}$.


\end{cor}





\begin{pf}


{\bf Необходимость.} Если оператор вложения


$\varphi^{*} :L_{p}^{1}(D^{\prime})\to L_{q}^{1}(D)$


ограничен, то $\varphi\in\ACL(D)$ и имеет конечное искажение.


Так же, как и в лемме 1, доказывается, что функция $\Xi$, определенная


перед следствием, является ограниченной счетно-аддитивной


функцией, определенной на открытых множествах $W\subset


D^{\prime}$. Кроме того, в силу очевидной оценки $\Phi(W)\leq


\Xi(W)$ и поточечного неравенства


$\alpha_{p,q}H^{\varkappa}_q(y)\leq\Phi'(y)$, установленного в


доказательстве теоремы 1, получим


$$


\alpha_{p,q}H^\varkappa_{p,q}(W)\leq \Phi(W)\leq \Xi(W)\leq


H^\varkappa_{p,q}(W).


$$


Отсюда вытекает абсолютная непрерывность


функции $\Xi$ и поточечная оценка


$$


\alpha^\varkappa_{p,q}


H^\varkappa_q(y)\leq \ov\Xi{}^{\prime}(y)\leq H^\varkappa_q(y)


$$


для почти всех $y\in D^{\prime}$. Далее, для произвольной


совокупности попарно непересекающихся открытых


множеств $W_i\subset A\subset W\subset D^{\prime}$ имеем


$$


\sum_{i=1}^\infty\Xi(W_i)\leq \sum_{i=1}^\infty\int_{W_i}


H^{\varkappa}_q(y)\,dy\leq


\int_{W}H^{\varkappa}_q(y)\,dy\leq \vartheta\Xi(W),


$$


где $\vartheta\geq1$ --- некоторая постоянная.


Следовательно, $\Xi$ --- абсолютно непрерывная счетно


$\vartheta$-квазиаддитивная функция множества.


\medskip





{\bf Достаточность.} Если дана функция множества, обладающая указанными


в формулировке следствия 2 свойствами, то в силу предложения 1


$$


\lambda_1H^\varkappa_{p,q}(D^{\prime})\leq


\int_{D^{\prime}}\ov\Xi{}^{\prime}(y)\,dy


\leq\Xi(D^{\prime})<\infty.


$$


Ограниченность оператора


$\varphi^{*} :L_{p}^{1}(D^{\prime})\to L_{q}^{1}(D)$


вытекает теперь из теоремы 1.


Из приведенных выше оценок вытекает также эквивалентность нормы оператора


$\varphi_W^*:L^1_p(W)\to L^1_q(\varphi^{-1}(W))$


выражению


$\Big(\int\limits_W\ov\Xi{}^{\prime}(y)\,dy\Big)^{\frac{1}{\varkappa}}$


для любого открытого множества $W\subset D^{\prime}$.


\end{pf}





Отображение $\varphi : D\to D^{\prime}$


называется {\it измеримым}, если прообраз измеримого множества измерим.





Для измеримого отображения $\varphi : D\to D^{\prime}$


определим {\it объемную производную} ``обратного отображения''


$$


J_{\varphi^{-1}}(y) = \lim_{r\to 0}


\frac{|\varphi^{-1}(B(y,r))|}{|B(y,r)|},\quad y\in D^{\prime}.


$$


Отображение $\varphi : D\to D^{\prime}$


{\it обладает ${\cal N}^{-1}$-свойством},


если прообраз всякого множества нулевой меры есть множество меры нуль.





Из ([28], теорема 4) и предложения 1, используя метод


доказательства теоремы 1, получаем 





%%%%% Предложение 4. ........








\input 02a-10.tex





\end{document}





